% 04_framework.tex — mirrors docs/06_Framework.md

\section{Framework Design}\label{sec:framework}

We design and implement a reproducible cross-tenant leakage
measurement framework with six cooperating modules plus a
reporter. The framework is the executable embodiment of the
threat model in \S\ref{sec:threat_model}: given a manifest of
tenants, attacks, and defenses, it spins up the reference
servers, drives concurrent tenant traffic, classifies
leakage events, and emits per-attack metrics.

\subsection{Goals}

\begin{itemize}
  \item \textbf{Replayable.} Given a manifest + seed, the
        framework produces identical results.
  \item \textbf{Tenant-aware.} First-class notion of
        \texttt{tenant\_id} across scheduler, evaluator, logger.
  \item \textbf{Multi-boundary.} Covers all eight MCP
        boundaries.
  \item \textbf{Pluggable.} New attacks, defenses, and metrics
        can be added without changing the core.
\end{itemize}

\subsection{Architecture}

\begin{table}[t]
\caption{Framework modules (\texttt{framework/}). Each module
has a single responsibility and a typed input/output contract.}
\label{tab:modules}
\centering
\small
\begin{tabularx}{\linewidth}{l X X}
\toprule
Module & Responsibility & File \\
\midrule
Scheduler    & Drive concurrent tenants + attack queue   & \texttt{scheduler/scheduler.py} \\
PayloadGen   & Deterministic per-seed payload mutation   & \texttt{scheduler/payloads.py} \\
Connector    & Speak MCP (stdio / HTTP+SSE / dummy)      & \texttt{target/connector.py} \\
Evaluator    & Classify leakage via substring + sha256   & \texttt{evaluator/evaluator.py} \\
Metrics      & Aggregate per boundary / defense          & \texttt{metrics/metrics.py} \\
Logger       & Persist every event as JSONL              & \texttt{logger/logger.py} \\
Reporter     & Render Markdown + HTML reports            & \texttt{reports/reporter.py} \\
\bottomrule
\end{tabularx}
\end{table}

A high-level sequence diagram is available in
\texttt{docs/diagrams/framework\_sequence.svg}. The diagram
covers the five prompt-required steps: framework spins up
Tenant~A and Tenant~B; payload injected via Tenant~A;
Tenant~B's state is sampled; oracle classifies result; logger
writes event.

\subsection{RunConfig schema}

The framework hydrates a \texttt{RunConfig} (Pydantic v2 model)
from a YAML manifest. The full schema is:

\begin{lstlisting}[language=yaml,caption={RunConfig schema (excerpt).},label=lst:runconfig]
run:
  seed: 42
  repeats: 30
  concurrency: 4
tenants:
  - id: tenant-A
    name: Alice
    allowed_tools: [echo]
  - id: tenant-B
    name: Bob
    allowed_tools: [echo]
attacks:
  - id: A-CCH-T
    parameters: {}
defenses:
  per_tenant_tool_registry: false
  tenant_prefixed_cache_keys: false
  resource_path_canonicalisation: false
  mtls: false
target:
  transport: dummy   # stdio | http_sse | streamable_http | dummy
output:
  log_dir: analysis/runs
  output_dir: analysis/runs
  log_format: csv
\end{lstlisting}

\texttt{RunConfig} enforces at least 2 tenants because the
harness requires a (source, sink) pair to detect leakage. The
\texttt{defenses} block is consumed by the reference
secure-server factory (see \S\ref{sec:evaluation}). The full
schema is versioned: \texttt{schema\_version: '1.0'} and
\texttt{dataset\_version: 'phase9-v1'} are mandatory fields;
the runner refuses to execute a manifest whose schema version
does not match the runner's expected version.\footnote{The
schema is documented in
\texttt{experiments/manifests/schema.yaml}. The pinning
protects against silent manifest-format changes between
Phases 9 and 12.} Sample size $n = 30$ per cell is the
pre-registered choice (\texttt{analysis/power.md}) that
targets medium effect sizes (Cliff's $\delta \geq 0.50$)
at power $\geq 0.80$ and $\alpha = 0.05$; smaller effects
are characterised as exploratory; the headline comparisons
in \S\ref{sec:evaluation} target $\delta \geq 0.80$ (large)
and are well within the $n = 30$ budget.

\subsection{Metrics}

The framework computes four metrics per attack per defense
level:

\begin{itemize}
  \item \textbf{\texttt{leakage\_rate}} -- fraction of call
        events in the bucket that have a paired leakage event
        (range $[0, 1]$).
  \item \textbf{\texttt{time\_to\_leak\_ms}} -- median latency
        between \texttt{attack\_started\_at} and
        \texttt{leakage\_detected\_at}; \texttt{null} if no
        leakage.
  \item \textbf{\texttt{defense\_overhead\_ms}} -- p95 of
        defended-call latency minus p95 of undefended-call
        latency; \texttt{null} if either side has no samples.
  \item \textbf{\texttt{utility\_retention}} -- $1 -
        \mathrm{errors\_defended}/\mathrm{errors\_undefended}$,
        clamped to $[0, 1]$.
\end{itemize}

\subsection{Reproducibility}

Each manifest declares a seed; the framework's payload
generator and the connector's probabilistic leak path both
consume that seed via \texttt{random.Random}. Defenses are
enabled/disabled at server startup, not at runtime, so the
server's behaviour is fully determined by
$(\mathrm{manifest}, \mathrm{seed})$. The artifact includes
pinned dependency versions (\texttt{pyproject.toml}) and a
\texttt{requirements.txt} for byte-identical reproduction.

\subsection{Reference servers}

We ship two reference servers:

\begin{itemize}
  \item \textbf{\texttt{mcp\_servers/vulnerable/}} --- a
        minimal MCP server with all eight boundaries
        deliberately mis-configured: shared tool registry,
        shared cache keys, no URI canonicalisation, no JWT
        \texttt{aud} claim, and no defense middleware.
  \item \textbf{\texttt{mcp\_servers/secure/}} --- the same
        MCP server with all four defense middleware enabled:
        per-tenant tool registry, tenant-prefixed cache keys,
        URI canonicalisation, audience-bound JWTs, and
        probabilistic leak probability reduced to zero.
\end{itemize}

The two servers share the same wire protocol, the same tool
catalogue, and the same fixtures; the only difference is the
defense middleware. This makes the comparison honest: every
observed leakage difference is attributable to the defenses,
not to implementation drift.

\paragraph{Tool catalogue.}
The reference servers expose eight tools (\texttt{echo},
\texttt{get\_secret}, \texttt{list\_tenants}, \texttt{set\_env},
\texttt{read\_file}, \texttt{write\_file}, \texttt{cache\_get},
\texttt{cache\_put}). The tool catalogue and per-tool
arguments are documented in
\texttt{artifact/release/TOOL\_CATALOGUE.md}; the
\texttt{list\_tenants} tool is admin-only on the secure
server.

\subsection{Why a measurable oracle, not a model-based judge}

A subtle design question is whether the Evaluator (oracle)
should be a learned model or a deterministic substring-and-hash
match. We chose the latter for three reasons:

\begin{enumerate}
  \item \textbf{Reproducibility.} A deterministic oracle
        produces bit-identical results across machines. A
        learned oracle's outputs depend on the model's
        checkpoint, which is a moving target.
  \item \textbf{Causal interpretability.} When the oracle
        flags a leakage event, the developer can inspect the
        matched excerpt and verify the leak path. A model-based
        oracle produces a probability that requires
        interpretation.
  \item \textbf{Speed.} The deterministic oracle runs in
        $< 1\,\mu s$ per event; a model-based oracle takes
        $\sim 50\,ms$ per event. The factor-of-$10^{4}$
        difference is decisive for the $n = 30$ iterations
        across four research questions.
\end{enumerate}

The trade-off is false negatives: a deterministic oracle
misses leakages whose payload does not match the seeded
marker. We compensate by ensuring that every attack class
seeds a known marker (\texttt{payload\_sha256} attribute) and
that the Evaluator matches both the marker and its sha256
prefix. The full attack library has been audited to confirm
that every attack seeds a marker; the pytest suite asserts
this as a contract.

\subsection{Concurrency model}

The Scheduler uses Python's \texttt{asyncio} with a bounded
\texttt{Semaphore} to drive concurrent tenant traffic. Each
attack iteration runs in its own task; the scheduler awaits
all tasks at the iteration boundary before flushing the
logger. The concurrency limit is set per-manifest
(\texttt{run.concurrency}); our experiments use
\texttt{concurrency = 4} for RQ-1, RQ-2, RQ-3, and
\texttt{concurrency = 8} for RQ-4 to stress-test the defense
composition.

The connector's leak-probability path is per-connection (not
global), so concurrent tenants do not share a seed. The
Evaluator's substring match is per-tenant-pair, so concurrent
tenants do not cross-contaminate. The Logger appends to a
single JSONL file with file-level locking; race conditions
are confined to write boundaries, which the iteration
flushes handle.

\subsection{What the framework does not do}

The framework is intentionally narrow:

\begin{itemize}
  \item It does not run a real LLM. The connector uses a
        deterministic in-process channel; prompt-injection
        attacks inject markers, not instructions to a model.
  \item It does not measure silicon-level side channels. The
        connector's timing is deterministic at the
        microsecond scale; Spectre/Meltdown-class attacks
        would require a separate harness.
  \item It does not fuzz the wire protocol beyond the
        \texttt{A-FUZZ-001} grammar attack. A full mutation
        campaign (à la LSPFuzz~\cite{zhu_2025_lspfuzz}) is
        future work.
  \item It does not handle multi-language SDKs. The reference
        connector is Python-only; the TypeScript and Rust
        SDKs are out of scope.
\end{itemize}

These exclusions are deliberate: they keep the framework
focused on the protocol-layer question and the empirical
results interpretable.